\begin{definition}[Матричный подсчёт маршрутов]
Обозначим \(w_k(i,j)=(A^k)_{ij}\). В отличие от простого пути, маршрут
допускает многократное посещение вершин и повторное прохождение рёбер.
Для кратных рёбер различается также выбор конкретного ребра.
\end{definition}
\begin{lemma}[Разбиение маршрута на две части]
Для \(p,q\ge0\)
\[
 w_{p+q}(i,j)=\sum_s w_p(i,s)w_q(s,j).
\]
\end{lemma}
\begin{proof}
Разрежем маршрут после \(p\)-го ребра. Вершина разреза \(v_s\) и обе
части определены однозначно. Обратно, любую пару таких частей можно
склеить. Для фиксированного \(s\) действует правило произведения,
затем правило суммы. Это комбинаторное доказательство \(A^{p+q}=A^pA^q\).
\end{proof}
\begin{graphillustration}{Каждый маршрут однозначно разбивается в $v_s$ на префикс длины $p$ и суффикс длины $q$; линии обозначают целые участки.}{fig:02-split}
\begin{tikzpicture}[font=\small]
\node[gi vertex] (i) at (0,0) {$v_i$};
\node[gi selected] (s) at (4,0) {$v_s$};
\node[gi vertex] (j) at (8,0) {$v_j$};
\draw[gi emphasis] (i) to[bend left=15] (s);
\draw[gi alternate] (s) to[bend left=15] (j);
\node[gi note] at (2,0.75) {$p$ рёбер};
\node[gi note] at (6,0.75) {$q$ рёбер};
\node[gi note] at (4,-0.8) {разрез после $p$-го ребра};
\end{tikzpicture}
\end{graphillustration}


\begin{theorem}[Полезные суммы]
Число маршрутов длины \(k\), начинающихся в \(v_i\), равно
\(\sum_j(A^k)_{ij}\); число всех маршрутов длины \(k\) равно
\(\boldsymbol1^{\mathsf T}A^k\boldsymbol1\); число замкнутых маршрутов
длины \(k\) с отмеченной начальной вершиной равно \(\operatorname{tr}A^k\).
\end{theorem}
\begin{proof}
В первом случае маршруты разбиты по последней вершине, во втором —
по упорядоченной паре концов. В третьем берутся только маршруты,
у которых концы совпадают, то есть диагональные элементы матрицы.
Основная теорема билета 1 превращает каждое слагаемое в соответствующее число маршрутов.
\end{proof}
\begin{cor}[Для простого неориентированного графа]
\[
 (A^2)_{ii}=d(v_i),\qquad (A^2)_{ij}=|N(v_i)\cap N(v_j)|\quad(i\ne j),
 \qquad \operatorname{tr}A^2=2m.
\]
Число треугольников равно \(\operatorname{tr}A^3/6\).
\end{cor}
\begin{proof}
Двухшаговый маршрут имеет вид \(v_i,v_s,v_j\); его середина — общий
сосед концов. При совпадении концов можно выбрать любого соседа.
Сумма степеней равна \(2m\). Замкнутый трёхшаговый маршрут в простом
графе проходит по трём различным вершинам треугольника; у каждого
треугольника три выбора начала и два направления обхода.
\end{proof}
\begin{example}[Треугольник]
Для \(K_3\)
\[
 A=\begin{pmatrix}0&1&1\\1&0&1\\1&1&0\end{pmatrix},\quad
 A^2=\begin{pmatrix}2&1&1\\1&2&1\\1&1&2\end{pmatrix},\quad
 A^3=\begin{pmatrix}2&3&3\\3&2&3\\3&3&2\end{pmatrix}.
\]
Из каждой вершины выходят четыре маршрута длины два; всего их 12.
След \(A^3\) равен 6, так что треугольник ровно один.
\end{example}
\begin{graphillustration}{Шесть замкнутых маршрутов длины 3 одного треугольника: три начала (двойной контур) и два направления.}{fig:02-triangle}
\begin{tikzpicture}[font=\small]
\begin{scope}[shift={(0.0,0.0)}]
\node[gi selected] (1) at (0,0) {$1$};
\node[gi vertex] (2) at (1,1.1) {$2$};
\node[gi vertex] (3) at (2,0) {$3$};
\draw[gi arrow] (1) -- (2);
\draw[gi arrow] (2) -- (3);
\draw[gi arrow] (3) -- (1);
\node[gi note] at (1,-0.65) {$1231$};
\end{scope}
\begin{scope}[shift={(4.2,0.0)}]
\node[gi vertex] (1) at (0,0) {$1$};
\node[gi selected] (2) at (1,1.1) {$2$};
\node[gi vertex] (3) at (2,0) {$3$};
\draw[gi arrow] (2) -- (3);
\draw[gi arrow] (3) -- (1);
\draw[gi arrow] (1) -- (2);
\node[gi note] at (1,-0.65) {$2312$};
\end{scope}
\begin{scope}[shift={(8.4,0.0)}]
\node[gi vertex] (1) at (0,0) {$1$};
\node[gi vertex] (2) at (1,1.1) {$2$};
\node[gi selected] (3) at (2,0) {$3$};
\draw[gi arrow] (3) -- (1);
\draw[gi arrow] (1) -- (2);
\draw[gi arrow] (2) -- (3);
\node[gi note] at (1,-0.65) {$3123$};
\end{scope}
\begin{scope}[shift={(0.0,-2.8)}]
\node[gi selected] (1) at (0,0) {$1$};
\node[gi vertex] (2) at (1,1.1) {$2$};
\node[gi vertex] (3) at (2,0) {$3$};
\draw[gi arrow] (1) -- (3);
\draw[gi arrow] (3) -- (2);
\draw[gi arrow] (2) -- (1);
\node[gi note] at (1,-0.65) {$1321$};
\end{scope}
\begin{scope}[shift={(4.2,-2.8)}]
\node[gi vertex] (1) at (0,0) {$1$};
\node[gi selected] (2) at (1,1.1) {$2$};
\node[gi vertex] (3) at (2,0) {$3$};
\draw[gi arrow] (2) -- (1);
\draw[gi arrow] (1) -- (3);
\draw[gi arrow] (3) -- (2);
\node[gi note] at (1,-0.65) {$2132$};
\end{scope}
\begin{scope}[shift={(8.4,-2.8)}]
\node[gi vertex] (1) at (0,0) {$1$};
\node[gi vertex] (2) at (1,1.1) {$2$};
\node[gi selected] (3) at (2,0) {$3$};
\draw[gi arrow] (3) -- (2);
\draw[gi arrow] (2) -- (1);
\draw[gi arrow] (1) -- (3);
\node[gi note] at (1,-0.65) {$3213$};
\end{scope}
\end{tikzpicture}
\end{graphillustration}


\textbf{Граница применимости.} Деление следа на \(2k\) при \(k\ge4\)
обычно не даёт число простых циклов: появляются замкнутые маршруты
с повторениями. Для орграфа формула для треугольников также меняется.
